\subsection{Марковский момент}

\begin{definition}
Пусть $T \subset \R$. Случайная величина $\tau$ со значениями в $T \cup \{+ \infty\}$ называется марковским моментом относительно фильтрации $\{\F_t\}_{t \in T}$, если имеет место включение $\{\tau \leq t\} \in \F_t \forall t$. Если $\tau < +\infty$, то $\tau$ называют моментом остановки.
\end{definition}

\begin{example}
Для процесса $\{\xi_n\}_{n \in \N}$ и порожденной им фильтрации для произвольного борелевского множества $B \subset \R$ положим $\tau_B(\omega) = min\{n: \xi_n(\omega) \in B\}$, если такие $n$ есть, иначе $\tau_B(\omega) = +\infty$.

Тогда $\tau_B$ -- марковский момент, ибо для всякого $t \in \N$ множество $\{\tau_B \leq t\}$ равно объединению множеств $\{\omega : \xi_n(\omega) \in B\}, n = 1, \dots, t$, входящих в $\F_{\leq t}$.
\end{example}

Если $\F_t$ интерпретировать как класс событий, о наступлении или ненаступлении которых известно к моменту времени $t$, то марковость можно понимать так: если при каком-то $\omega$ марковский момент принял значение $t$, то к моменту времени $t$ это уже известно. В случае фильтрации, порожденной случайным процессом $\{\xi_t\}$, наглядный смысл таков: по наблюдениям значений процесса до момента времени $t$ включительно можно понять, наступил ли уже момент $\tau$ или попало ли $\tau$ в множество $B$.

\subsection{Теорема об остановке}

\begin{theorem}[об остановке]
Пусть $\{\xi_n\}_{n \in \Z_+}$ -- мартингал относительно фильтрации $\{\F_n\}_{n \in \Z_+}$ и $\tau$ -- ограниченный момент остановки. Тогда имеем $E \xi_{\tau} = E \xi_0$. В частности, если $E \xi_0 = 0$, то $E \xi_{\tau} = 0$.
\end{theorem}

\begin{proof}
Пусть $\tau \leq m$. Тогда $|\xi_{\tau}| \leq |\xi_1| + \dots + |\xi_m|$, поэтому $E \xi_{\tau}$ существует. Для фиксированного $j \leq m$ имеем

$$E(\xi_{\tau}I(\{\tau = j\})) = E(\xi_{j}I(\{\tau = j\})) = E(\xi_{j + 1}I(\{\tau = j\})) = \dots = E(\xi_{m}I(\{\tau = j\}))$$, поскольку $\{\tau = j\} \in \F_{j}$. Сумма по $j$ от $0$ до $m$ дает слева $E \xi_{\tau}$, а справа $E \xi_m = E \xi_0$
\end{proof}

\begin{definition}
    Пусть $\{\F_t\}_{t \geq 0}$ -- фильтрация. Отображение $\tau \to [0, +\infty]$ называется опциональным моментом, если $\{\tau < t\} \subset \F_t$
\end{definition}

\begin{example}
    Пусть процесс $\{\xi_t\}_{t \geq 0}$ согласован с фильтрацией $\{\F_t\}_{t \geq 0}$ и имеет непрерывные справа траектории. Тогда для открытого множество $U$ момент $\tau_U = inf\{t : \xi_t \in U\}$ является опциональным. Действительно, $\tau_U < t$ в точности тогда, когда есть $s < t$ с $\xi_s \in U$, причем из-за непрерывности справа такое $s$ можно взять рациональным. Поэтому $\tau_U < t$ есть счетное объединение множеств $\{\xi_s \in U\} \in \F_s$ по рациональным $s < t$.
\end{example}

\begin{theorem}
    Пусть $\{\xi_t\}_{t \geq 0}$ -- мартингал относительно фильтрации $\{\F_t\}_{t \geq 0}$, имеющий непрерывные справа траектории. Тогда для всякого ограниченного опционального момента имеем $E \xi_{\tau} = E \xi_0$
\end{theorem}